% TOP_PROBLEM: {"id": 2647, "title": "Kourovka Notebook Problem 21.138", "category_id": 17, "category": {"id": 17, "name": "group_theory", "display_name": "Group Theory", "description": "Problems about groups, group actions, representations, and related algebraic structures.", "slug": "group-theory", "order_index": 17, "created_at": "2026-07-31T15:26:25.671Z"}, "status": "open", "problem_number": "KOU-21.138", "set_id": 10}
% FIRST_POSTED: 2026-10-01
% ENTRY_KIND: statement_only
% UnsolvedMath ID 2647; source catalogue code KOU-21.138
% !TeX program = lualatex
% Definition and sources only; this file is not an LLM solution attempt.
\documentclass[11pt,a4paper]{article}
\usepackage[margin=25mm]{geometry}
\usepackage{fontspec}
\setmainfont{lmroman10-regular.otf}[BoldFont=lmroman10-bold.otf,
 ItalicFont=lmroman10-italic.otf,BoldItalicFont=lmroman10-bolditalic.otf]
\usepackage{amsmath,amssymb,mathtools}
\usepackage{unicode-math}
\setmathfont{latinmodern-math.otf}
\AtBeginDocument{\let\setminus\smallsetminus}
\usepackage{xurl,hyperref}
\hypersetup{colorlinks=true,urlcolor=blue}
\setcounter{secnumdepth}{0}
\setlength{\parindent}{0pt}
\setlength{\parskip}{5pt}
\setlength{\emergencystretch}{3em}
\begin{document}
\section*{Kourovka Notebook Problem 21.138}\label{2647}
{\small UnsolvedMath ID 2647 · KOU-21.138 · Group Theory · Kourovka Notebook - New Problems, Issue 21}
\subsection{Definitions and mathematical statement}
A group is finitely presented if it has a presentation $\langle S\mid R\rangle$ with both the generating set $S$ and the list of defining relators $R$ finite. For a subgroup $H\le G$, its index $[G:H]$ is the cardinality of the set of left cosets. Infinite index means that this coset set is infinite.

A free group on a set $B$ is a group in which arbitrary assignments of $B$ into any group extend uniquely to homomorphisms. Free subgroups in the hypothesis may have any rank, including zero, finite rank, or infinite rank. A surface group here means the fundamental group of a closed connected aspherical surface, orientable or nonorientable. Aspherical means its universal cover is contractible; thus this convention includes the torus and Klein bottle and excludes the sphere and projective plane.

Suppose $G$ is infinite and finitely presented, and every subgroup $H\le G$ of infinite index is a free group. Must $G$ itself be either a free group or a surface group?

The hypothesis is imposed on all infinite-index subgroups, not just finitely generated ones. Finite-index subgroups are not required to be free. The desired conclusion is an isomorphism of $G$ itself with one of the two indicated kinds of groups, rather than just existence of a finite-index subgroup of such a kind. No one-relator, residual-finiteness, or geometric-action assumption is present.
\subsection{Short English statement}
Are free groups and closed aspherical surface groups the only infinite finitely presented groups whose every infinite-index subgroup is free?
\subsection{Sources}
\begin{itemize}
\item[S1] ULAM.AI and the UnsolvedMath Contributors, supplied problems.json, record 2647 (KOU-21.138), Kourovka Notebook - New Problems, Issue 21. Catalogue identity and source transcription; CC BY 4.0.
\url{https://huggingface.co/datasets/ulamai/UnsolvedMath}.
\item[S2] The Kourovka Notebook, 21st edition, new problems, Problem 21.138. Expanded formulation of the supplied catalogue statement.
\url{https://arxiv.org/pdf/1401.0300v47}.
\item[S3] H. Wilton, Surface groups among cubulated hyperbolic and one-relator groups, Question 0.1; specifies closed aspherical surfaces and distinguishes known restricted cases.
\url{https://arxiv.org/html/2406.02121v3}.
\end{itemize}
\end{document}
